% =====================================================================
%  Capitolo 1 — blocchi da sostituire e da inserire
%
%  Non e' un capitolo nuovo: sono i quattro interventi puntuali sul
%  Cap1.tex esistente. Ogni blocco riporta il testo attuale in commento
%  e la versione proposta sotto.
% =====================================================================


% ---------------------------------------------------------------------
% INTERVENTO 1 — nuova sottosezione, da inserire alla fine di
% "Sfide in scenari reali", prima di "Obiettivi".
%
% Motivo: il capitolo elenca domain shift e occlusioni, ma non il
% problema che in pratica fa fallire i sistemi installati e che il
% capitolo sperimentale misura.
% ---------------------------------------------------------------------

\subsection{La prevalenza dell'evento e il costo dei falsi allarmi}

Un terzo problema, meno discusso in letteratura ma decisivo nella pratica,
riguarda la frequenza dell'evento da rilevare.

Nei dataset di riferimento le cadute sono abbondanti per costruzione: in
UP-Fall, che è l'archivio utilizzato in questo lavoro, le finestre temporali
contenenti una caduta o una condizione post-caduta sono circa il 12\% del
totale. È una proporzione necessaria per addestrare un classificatore, ma non ha
alcun rapporto con la realtà di un'abitazione, dove una persona anziana cade
mediamente una o due volte l'anno.

Lo squilibrio si traduce in un vincolo progettuale stringente. Un sistema che
produce una decisione ogni frazione di secondo compie, nell'arco di un anno,
alcune decine di milioni di valutazioni, delle quali soltanto qualche decina
corrisponde a un evento reale. A parità di specificità, lo stesso modello che
sul dataset appare accurato può quindi produrre in ambiente domestico un numero
di falsi allarmi di diversi ordini di grandezza superiore a quello tollerabile.

Il riferimento pratico è severo. Gli studi longitudinali su installazioni
funzionanti riportano tassi dell'ordine di $0{,}05$ eventi di falso allarme per
ora di utilizzo \cite{bagala2012, palmerini2015}, e il criterio clinicamente
accettabile è di uno o due episodi alla settimana. Un sistema che superi
sistematicamente questa soglia viene disattivato dall'utente, e la sua
sensibilità nominale diventa irrilevante.

Ne segue che la valutazione di un sistema di rilevamento non può limitarsi alle
metriche di classificazione calcolate sul dataset. Occorre distinguere fra le
\emph{prestazioni del classificatore}, misurate per finestra temporale, e le
\emph{prestazioni del sistema}, misurate in termini di eventi di allarme per
unità di tempo — grandezze che, come mostrato nel
Capitolo~\ref{cap:sperimentale}, differiscono fra loro di quasi due ordini di
grandezza.


% ---------------------------------------------------------------------
% INTERVENTO 2 — sostituzione del secondo punto dell'elenco in
% "Obiettivi".
%
% TESTO ATTUALE (da rimuovere):
%   \item \textbf{Classificazione temporale tramite rete leggera}:
%   analizzare le sequenze di feature tramite una rete ricorrente GRU,
%   scelta come compromesso tra capacita di modellare dipendenze
%   temporali e costo computazionale contenuto.
%
% Motivo: e' formulato come una scelta di progetto conclusa, mentre il
% capitolo sperimentale mostra che l'architettura non e' distinguibile
% da nessun'altra e che il sistema installato usa un classificatore
% diverso. Il lettore troverebbe smentita una promessa del Capitolo 1.
% ---------------------------------------------------------------------

\item \textbf{Classificazione temporale tramite rete leggera}: analizzare le
sequenze di grandezze tramite una rete ricorrente GRU, adottata come
\emph{ipotesi di lavoro}. L'ipotesi è che una rete ricorrente compatta sia in
grado di modellare le dipendenze temporali dell'evento a un costo compatibile
con l'esecuzione su dispositivo. La verifica sperimentale di questa ipotesi, e
il suo esito, sono oggetto del Capitolo~\ref{cap:sperimentale}.


% ---------------------------------------------------------------------
% INTERVENTO 3 — integrazione del primo punto dei requisiti in
% "Requisiti di un sistema automatico di fall detection".
%
% TESTO ATTUALE (da sostituire):
%   \item In primo luogo deve essere in grado di rilevare in maniera
%   tempestiva gli eventi critici, generando alert quando viene
%   individuata una caduta.
%
% Motivo: il requisito e' qualitativo e non fornisce un metro con cui
% il capitolo sperimentale possa misurarsi.
% ---------------------------------------------------------------------

\item In primo luogo deve rilevare tempestivamente gli eventi critici,
generando un allarme quando viene individuata una caduta, e deve farlo senza
produrre un numero di segnalazioni improprie tale da comprometterne
l'accettabilità. I due requisiti sono in tensione e vanno dichiarati insieme:
in un contesto \emph{safety-critical} un mancato rilevamento è l'errore più
grave, ma un sistema che segnali ripetutamente eventi inesistenti viene
disattivato, e la sua sensibilità nominale perde ogni valore. Il riferimento
adottato in questo lavoro è quello degli studi su installazioni reali
funzionanti, dell'ordine di $0{,}05$ eventi di falso allarme per ora.


% ---------------------------------------------------------------------
% INTERVENTO 4 — aggiunta alla fine di "Conseguenze cliniche ed
% esistenziali", dopo il paragrafo sulla long-lie syndrome.
%
% Motivo: la long-lie syndrome e' gia' citata, ma il capitolo non trae
% la conseguenza progettuale che il capitolo sperimentale utilizza.
% Piantare qui il seme rende leggibile la struttura a due vie descritta
% nelle conclusioni.
% ---------------------------------------------------------------------

Questa osservazione ha una conseguenza progettuale che sarà ripresa nel seguito.
L'evento da rilevare si compone infatti di due fasi distinte: la \emph{caduta}
propriamente detta, che dura pochi secondi, e la \emph{permanenza a terra}, che
può durare ore ed è la fase clinicamente determinante. Le due fasi hanno
proprietà molto diverse dal punto di vista del riconoscimento automatico — la
prima è una transizione dinamica, la seconda una condizione statica — e la
scelta di quale delle due il sistema debba riconoscere non è neutra. Il
Capitolo~\ref{cap:sperimentale} mostra che questa scelta produce sull'esito
finale un effetto superiore a quello di qualunque altra variabile considerata in
questo lavoro.
